% 原书第六篇6.4节：分子与晶体167--180题。
% 底本 Irodov Problems in General Physics 印刷页251--253，PDF页260--262；
% 原书书末答案印刷346--347，PDF355--356。
\moleculeproblem{167}{251}
双硫分子 $\mathrm{S_2}$ 的两核距离 $d=189\,\mathrm{pm}$。
求其被激发到第一个转动激发态时的角速度。
\exsolution
两个相同硫原子质量为 $m_S$，绕质心转动惯量
$I=2m_S(d/2)^2=m_Sd^2/2$。
量子转子角动量模为
$L_J=\hbar\sqrt{J(J+1)}$；第一个激发态 $J=1$，
所以 $L_1=\sqrt2\hbar$，而 $L=I\Omega$。
\[
 \boxed{\Omega=\frac{\sqrt2\hbar}{I}
 =\frac{2\sqrt2\hbar}{m_Sd^2}
 \simeq1.57\times10^{11}\,\mathrm{s^{-1}}}.
\]
\excheck 使用角动量的模 $\sqrt{J(J+1)}\hbar$，
不能在本题直接把 $J\hbar$ 视为转动角动量。

\moleculeproblem{168}{251}
HCl 分子的相邻两个转动能级之差为
$\Delta E=7.86\,\mathrm{meV}$，核间距
$d=127.5\,\mathrm{pm}$。
求两能级的转动量子数。
\exsolution
刚性转子的能级为
$E_J=\hbar^2J(J+1)/(2I)$，
$I=\mu d^2$，
$\mu=m_Hm_{\mathrm{Cl}}/(m_H+m_{\mathrm{Cl}})$。
于是
\[
 E_{J+1}-E_J=\frac{\hbar^2}{I}(J+1),\qquad
 J+1=\frac{I\Delta E}{\hbar^2}.
\]
代入 $m_H\simeq1.008u$、$m_{\mathrm{Cl}}\simeq35.45u$
得 $J+1\simeq3.00$，故
\[
 \boxed{J=2,\quad J'=3}.
\]
\excheck 能级间隔与 $J+1$ 成正比，非恒定能量间隔。

\moleculeproblem{169}{251}
氧分子 $\mathrm{O_2}$ 核间距为
$d=121\,\mathrm{pm}$，转动能量
$E_{\mathrm{rot}}=2.16\,\mathrm{meV}$。
求转动角动量模。
\exsolution
两氧核等质量，转动惯量
$I=m_Od^2/2$。
由经典转动能与量子转子能级一致的关系
$E=L^2/(2I)$，有
\[
 \boxed{L=\sqrt{2IE}
 =\sqrt{m_Od^2E}
 \simeq3.48\hbar\approx3.5\hbar}.
\]
\excheck 此处 $L$ 是角动量的模，
单位 $\mathrm{J\,s}$；不能仅凭能量直接令 $L=J\hbar$。

\moleculeproblem{170}{251}
证明双原子分子的纯转动谱相邻谱线角频率间隔相等。
若 CH 分子的谱线角频率间隔
$\Delta\omega=5.47\times10^{12}\,\mathrm{s^{-1}}$，
求转动惯量与核间距。
\exsolution
转动跃迁遵守 $\Delta J=1$，
其光子角频率
\[
 \omega_{J\to J+1}=
 \frac{E_{J+1}-E_J}{\hbar}
 =\frac{\hbar}{I}(J+1).
\]
因此相邻谱线的间隔
$\Delta\omega=\hbar/I$，
与 $J$ 无关。由此
\[
 \boxed{I=\frac{\hbar}{\Delta\omega}
 \simeq1.93\times10^{-47}\,\mathrm{kg\,m^2}
 =1.93\times10^{-40}\,\mathrm{g\,cm^2}}.
\]
设约化质量 $\mu=m_Cm_H/(m_C+m_H)$，
$I=\mu d^2$，从而
\[
 \boxed{d=\sqrt{I/\mu}\simeq112\,\mathrm{pm}}.
\]
\excheck 角频率差对应 $\hbar$ 而非普通频率差对应的 $h$。

\moleculeproblem{171}{252}
HF 分子的自然振动角频率
$\omega_v=7.79\times10^{14}\,\mathrm{s^{-1}}$，
核间距 $d=91.7\,\mathrm{pm}$。
忽略转动--振动耦合，
问从振动基态到第一振动激发态之间
容纳多少个转动激发能级？
\exsolution
第一振动能隙为 $\hbar\omega_v$。
转动能量
$E_J-E_0=\hbar^2J(J+1)/(2I)$，
$I=\mu d^2$，$\mu=m_Hm_F/(m_H+m_F)$。
低于振动激发能的转动能级满足
\[
 J(J+1)<\frac{2I\omega_v}{\hbar}.
\]
代入数据解得
$J<13.56$，允许正整数 $J=1,2,\ldots,13$，
因此
\[
 \boxed{13\text{ 个转动激发能级}}.
\]
\excheck 原有 $J=0$ 是转动基态，
不计入题目要求的“在其间的激发能级”。

\moleculeproblem{172}{252}
CO 分子自然振动角频率
$\omega_v=4.09\times10^{14}\,\mathrm{s^{-1}}$，
转动惯量 $I=1.44\times10^{-39}\,\mathrm{g\,cm^2}$。
估算纯转动谱能够出现多少条谱线。
\exsolution
约束转动能不超过一个振动量子
$\hbar\omega_v$，即
\[
 \frac{\hbar^2}{2I}J(J+1)<\hbar\omega_v.
\]
对于 $J\gg1$，得到
$J_{\max}\simeq\sqrt{2I\omega_v/\hbar}$。
纯转动跃迁由相邻 $J$ 能级产生，
因此光谱线数约为
\[
 \boxed{N\simeq\sqrt{\frac{2I\omega_v}{\hbar}}
 \simeq33}.
\]
\excheck 估计线数需要对
$J(J+1)$ 开平方；若误用 $2I\omega_v/\hbar$
直接计数，会高估到千余条。

\moleculeproblem{173}{252}
推导双原子分子转动能量 $E$ 处
单位能量间隔中的转动能级数 $dN/dE$，
并计算核间距 $d=267\,\mathrm{pm}$、
处于 $J=10$ 的 $\mathrm{I_2}$ 分子该值。
\exsolution
能量 $E_J=\hbar^2J(J+1)/(2I)$，
高 $J$ 下将 $J$ 视为连续变量：
\[
 \frac{dN}{dE}\simeq\frac{dJ}{dE}
 =\frac{2I}{\hbar^2(2J+1)}
 =\frac{2I}{\hbar^2\sqrt{1+8IE/\hbar^2}}.
\]
若 $J\gg1$，进一步化为
$\sqrt{I/(2\hbar^2E)}$。
对 $\mathrm{I_2}$，$I=m_Id^2/2$；
$J=10$ 时
\[
 \boxed{\frac{dN}{dE}\simeq
 \frac{I}{10.5\hbar^2}
 \simeq1.03\times10^4\,\mathrm{eV^{-1}}}.
\]
\excheck 独立反求 $J(E)=[\sqrt{1+8IE/\hbar^2}-1]/2$ 并微分，
得到 $dJ/dE=2I/[\hbar^2\sqrt{1+8IE/\hbar^2}]$；
此前把根号形式的分子误写为 $I$，漏掉因子2，现已纠正。
这里的能级密度不计每个 $J$ 的
$2J+1$ 磁量子态简并度；若计态密度须另乘简并度。

\moleculeproblem{174}{252}
求将双原子分子激发到第一振动激发态
和第一转动激发态所需能量之比。
分别估算 $\mathrm{H_2}$（$\omega=8.3\times10^{14}\,\mathrm{s^{-1}}$、
$d=74\,\mathrm{pm}$）、HI（$4.35\times10^{14}$、
$160\,\mathrm{pm}$）、$\mathrm{I_2}$（$0.40\times10^{14}$、
$267\,\mathrm{pm}$）。
\exsolution
第一振动激发能 $\Delta E_v=\hbar\omega$；
第一转动激发能为
$E_1-E_0=\hbar^2/I$，其中
$I=\mu d^2$。所以
\[
 \boxed{\frac{\Delta E_v}{\Delta E_r}
 =\frac{I\omega}{\hbar}
 =\frac{\mu d^2\omega}{\hbar}}.
\]
取对应约化质量代入，
\[
 \boxed{\mathrm{H_2}:36,\qquad
 \mathrm{HI}:1.75\times10^2,\qquad
 \mathrm{I_2}:2.85\times10^3}.
\]
\excheck 振动量子通常显著高于转动量子；
原书 $\mathrm{I_2}$ 给约 $2.9\times10^3$。

\moleculeproblem{175}{252}
氢分子振动角频率
$\omega=8.25\times10^{14}\,\mathrm{s^{-1}}$，
核间距 $d=74\,\mathrm{pm}$。
在温度 $T=875\,\mathrm K$ 时，
求处于第一振动激发能级 $v=1$
与第一转动激发能级 $J=1$ 的分子数比。
转动简并度为 $2J+1$。
\exsolution
振动 $v=1$ 与振动基态之差为
$\hbar\omega$；
第一转动激发能为
$\hbar^2/I=2B$，记 $B=\hbar^2/(2I)$。
忽略两自由度耦合，取相应激发态热平衡比例，
转动态 $J=1$ 的权重为3，
故
\[
 \boxed{\frac{N_{v=1}}{N_{J=1}}
 =\frac13\exp\!\left[
 -\frac{\hbar\omega-\hbar^2/I}{k_BT}\right]
 \simeq3.0\times10^{-4}}.
\]
\excheck 必须纳入 $J=1$ 的三重转动简并，
否则比例高估3倍。

\moleculeproblem{176}{252}
用玻尔兹曼分布推导谐振子的平均能量，
并求双原子气体摩尔振动热容 $C_{V,\mathrm{vib}}$。
对氯气 $\mathrm{Cl_2}$，$\omega=1.064\times10^{14}\,\mathrm{s^{-1}}$，
$T=300\,\mathrm K$，计算数值。
\exsolution
设 $\beta=1/(k_BT)$，
$E_v=\hbar\omega(v+1/2)$，
配分函数
\[
 Z=\sum_{v=0}^\infty e^{-\beta E_v}
 =\frac{e^{-\beta\hbar\omega/2}}{1-e^{-\beta\hbar\omega}}.
\]
平均能量
$\overline E=-\partial_\beta\ln Z$，得
\[
 \boxed{\overline E=
 \frac{\hbar\omega}{2}
 +\frac{\hbar\omega}{e^{\hbar\omega/(k_BT)}-1}}.
\]
取 $x=\hbar\omega/(k_BT)$，
对一摩尔的 $N_A\overline E$ 求导，
\[
 \boxed{C_{V,\mathrm{vib}}=
 R\frac{x^2e^x}{(e^x-1)^2}}.
\]
氯气 $x\simeq2.709$，所以
$\boxed{C_{V,\mathrm{vib}}\simeq0.561R
 \simeq4.66\,\mathrm{J/(mol\,K)}}$。
\excheck 零点能 $\hbar\omega/2$ 虽计入内能，
但与温度无关，对热容没有贡献。

\moleculeproblem{177}{252}
HCl 发射谱转动--振动带中心的零级线不允许出现，
相邻谱线角频率间隔
$\Delta\omega=0.79\times10^{13}\,\mathrm{s^{-1}}$。
求 HCl 核间距。
\exsolution
振转跃迁遵守 $\Delta J=\pm1$，
两支谱线在带中心邻线的角频率间隔等于
$2\hbar/I$，其中 $I=\mu d^2$。
因此
\[
 \boxed{d=\sqrt{\frac{2\hbar}{\mu\Delta\omega}}
 \simeq1.28\times10^{-10}\,\mathrm m
 \simeq0.13\,\mathrm{nm}}.
\]
\excheck 使用带中心缺失零级线的几何事实
来辨认给定间隔；换用普通频率需乘 $2\pi$。

\moleculeproblem{178}{252}
单色入射光波长 $\lambda_0=404.7\,\mathrm{nm}$，
氟分子 $\mathrm{F_2}$ 的自然振动角频率
$\omega_v=2.15\times10^{14}\,\mathrm{s^{-1}}$。
求最靠近未移频线的红、紫拉曼伴线波长。
\exsolution
斯托克斯伴线光子频率减少振动频率，
反斯托克斯线频率增加振动频率，
\[
 \frac{2\pi c}{\lambda_r}=
 \frac{2\pi c}{\lambda_0}-\omega_v,\qquad
 \frac{2\pi c}{\lambda_v}=
 \frac{2\pi c}{\lambda_0}+\omega_v.
\]
解得
\[
 \boxed{\lambda_r=
 \frac{\lambda_0}{1-\omega_v\lambda_0/(2\pi c)}
 \simeq424\,\mathrm{nm}},\quad
 \boxed{\lambda_v=
 \frac{\lambda_0}{1+\omega_v\lambda_0/(2\pi c)}
 \simeq387\,\mathrm{nm}}.
\]
\excheck 红伴线的光子能量低、波长长；
紫伴线相反，符号不可交换。

\moleculeproblem{179}{253}
双硫分子 $\mathrm{S_2}$ 的拉曼红、紫近邻伴线波长
分别为 $346.6\,\mathrm{nm}$ 和 $330.0\,\mathrm{nm}$。
求分子的自然振动角频率与拟弹性回复力系数。
\exsolution
拉曼两线对入射角频率对称，
两者角频率之差为 $2\omega_v$，于是
\[
 \boxed{\omega_v=\pi c\left(
 \frac1{\lambda_v}-\frac1{\lambda_r}\right)
 \simeq1.367\times10^{14}\,\mathrm{s^{-1}}}.
\]
近似谐振子有
$\omega_v^2=\kappa/\mu$，
两同质量硫原子的约化质量 $\mu=m_S/2$，
故
\[
 \boxed{\kappa=\mu\omega_v^2
 \simeq4.97\times10^2\,\mathrm{N/m}
 =4.97\,\mathrm{N/cm}}.
\]
\excheck 将 $\mathrm{N/m}$ 换成 $\mathrm{N/cm}$
必须除以100，不可漏换单位。

\moleculeproblem{180}{253}
对 $\mathrm{Cl_2}$ 分子的拉曼近邻紫、红伴线，
在 $T=300\,\mathrm K$、
$\omega_v=1.06\times10^{14}\,\mathrm{s^{-1}}$ 时，
求强度比 $I_v/I_r$，
以及温度加倍后该比值增加倍数。
\exsolution
斯托克斯（红线）强度与
$v\to v+1$ 的占据权重成正比；
反斯托克斯（紫线）强度与
$v+1\to v$ 的占据权重成正比。
谐振子热平衡给
$N_{v+1}/N_v=e^{-\hbar\omega_v/(k_BT)}$，
因而
\[
 \boxed{\frac{I_v}{I_r}
 =e^{-\hbar\omega_v/(k_BT)}
 \simeq0.0673}.
\]
温度倍增到 $2T$ 时，
\[
 \boxed{\frac{(I_v/I_r)_{2T}}{(I_v/I_r)_T}
 =e^{\hbar\omega_v/(2k_BT)}\simeq3.86}.
\]
\excheck 采用原书近似，忽略两伴线光频
四次幂的微小散射效率修正。
